\chapter{Cord Blood Testing and Banking}
\section*{What Is This Test?}
Cord blood testing examines blood collected from the umbilical cord and placenta after birth. It can be tested for blood-group or blood-cell findings, or collected and stored as a source of blood-forming stem cells.
\section*{Alternative Names}
\begin{itemize}\item Umbilical Cord Blood Testing\item Cord Blood Banking\item Hematopoietic Stem Cell Collection\end{itemize}
\section*{Principle}
Blood remaining in the cord and placenta is collected into a sterile bag after delivery. Testing may include cell counts, blood group, infectious-disease screening, HLA typing, and cell-dose or viability measurements for banking.
\section*{Why This Test Is Ordered}
Testing may evaluate a newborn's anemia, jaundice, blood-group incompatibility, or suspected infection. Banking may be considered when a family has a medical indication for a related donor or chooses public donation where available.
\section*{Primary vs Secondary Test}
Cord blood results are secondary to the newborn's examination and, when needed, a peripheral neonatal blood sample. Banking is a collection and storage service rather than a diagnostic test for the infant.
\section*{How This Test Is Performed}
After the cord is clamped and cut, trained staff draw blood from the cord vessels without entering the mother or baby. The unit is labeled, transported, tested, processed, and donated or stored under controlled conditions.
\section*{What Preparation Is Needed}
Families considering banking should review consent, collection logistics, costs, release policies, accreditation, and testing requirements before labor. Collection must not compromise routine care.
\section*{Contraindications and Precautions}
Collection may fail when volume is too small or urgent care takes priority. Banking does not guarantee that a unit will be suitable or sufficient for a future transplant.
\section*{Understanding Results}
Reports may include blood group, CBC, infectious markers, HLA type, total nucleated cells, CD34-positive cell count, viability, and volume. Interpretation depends on intended use.
\section*{Follow-up Tests}
Follow-up may include newborn blood testing, bilirubin and hemolysis studies, confirmatory infectious testing, HLA typing, or transplant-center evaluation.
\section*{References}
\begin{enumerate}\item MedlinePlus. Cord Blood Testing and Banking.\item Food and Drug Administration. Cord blood banking and stem-cell products.\end{enumerate}
\clearpage\section*{Understanding Results and Follow-up}
The laboratory report should identify the specimen, method, units, reference interval or decision limit, and whether the result is positive, negative, abnormal, or inconclusive. Reference intervals vary with age, sex, pregnancy, specimen type, assay, and laboratory; a result outside a range is not automatically a diagnosis, and a result within a range does not always exclude disease. Discuss unexpected results with the ordering clinician, who can decide whether repeat or confirmatory testing, treatment, or referral is appropriate. Do not change medicines or delay urgent care based on an isolated result.

\section*{Regulatory Status and Authorization Date}This chapter describes a test category, not a specific commercial product. FDA, CDSCO, EU IVDR, and MHRA approval or registration dates therefore cannot be assigned without the assay manufacturer, product name, instrument, and intended use. For laboratory-developed tests, an individual agency approval date may not exist. See the Preface for the interpretation of regulatory status.
\clearpage
